\begin{lemma}
\label{editorial-source-lemma-cover-module}
Let $R$ be a ring. Let $f_1, \ldots, f_n$ be elements of $R$
generating the unit ideal. Let $M$ be an $R$-module.
The sequence
$$
0 \to
M \xrightarrow{\alpha}
\bigoplus\nolimits_{i = 1}^n M_{f_i} \xrightarrow{\beta}
\bigoplus\nolimits_{i, j = 1}^n M_{f_i f_j}
$$
is exact, where $\alpha(m) = (m/1, \ldots, m/1)$
and $\beta(m_1/f_1^{e_1}, \ldots, m_n/f_n^{e_n})
= (m_i/f_i^{e_i} - m_j/f_j^{e_j})_{(i, j)}$.
\end{lemma}

\begin{proof}
It suffices to prove exactness after localization at each maximal
ideal $\mathfrak m$, by Lemma \ref{algebra-lemma-characterize-zero-local}.
The finite direct sums commute with localization. For each original
index $i$, the comparison of the two localization orders is
$$
(M_{f_i})_{\mathfrak m}\longrightarrow
(M_{\mathfrak m})_{f_i/1},\qquad
\frac{x/f_i^a}{s}\longmapsto\frac{x/s}{(f_i/1)^a},
$$
where $s\in R\setminus\mathfrak m$. Its reverse formula is the
inverse, by Proposition \ref{editorial-source-proposition-localize-twice-module}.
The same formulas with $f_if_j$ give the comparisons on the last term.
They commute with the original maps $\alpha$ and $\beta$, including
the difference in each ordered pair $(i,j)$.

\medskip\noindent
Choose an original index $k$ with $f_k\notin\mathfrak m$, which exists
because $(f_1,\ldots,f_n)=R$. Keep all indices and all elements $f_i$.
The natural map
$\iota_k:M_{\mathfrak m}\to(M_{f_k})_{\mathfrak m}$
has the explicit inverse
$$
\theta_k\left(\frac{x/f_k^a}{s}\right)=\frac{x}{sf_k^a}
\quad\text{in }M_{\mathfrak m}.
$$
The denominator $sf_k^a$ lies outside $\mathfrak m$.
For each $i$, the natural map
$(M_{f_i})_{\mathfrak m}\to(M_{f_kf_i})_{\mathfrak m}$
is also an isomorphism, with inverse
$$
\theta_{k,i}\left(\frac{x/(f_kf_i)^a}{s}\right)
=\frac{x/f_i^a}{sf_k^a}.
$$
Its forward formula sends $(x/f_i^a)/s$ to
$(f_k^ax/(f_kf_i)^a)/s$. The two composites are identities by
the equality-of-fractions relations. These maps come from localization
at elements which already act invertibly, so the formulas are
well-defined and preserve the original module structures.

\medskip\noindent
Projection to the $k$th summand followed by $\theta_k$ is a left
inverse of $\alpha_{\mathfrak m}$, proving injectivity.
Now let $(z_i)_{i=1}^n$ be in the kernel of $\beta_{\mathfrak m}$
and put $z=\theta_k(z_k)\in M_{\mathfrak m}$.
The $(k,i)$ coordinate says that the images of $z_k$ and $z_i$
in $(M_{f_kf_i})_{\mathfrak m}$ agree. Applying $\theta_{k,i}$
shows that $z_i$ is the image of $z$ in $(M_{f_i})_{\mathfrak m}$.
For completeness, on a representative $z_k=(x/f_k^a)/s$ the image
in the pair localization is $(f_i^ax/(f_kf_i)^a)/s$; applying
$\theta_{k,i}$ gives $(f_i^ax/f_i^a)/(sf_k^a)$, which is precisely
the image of $x/(sf_k^a)=z$. Thus
$\alpha_{\mathfrak m}(z)=(z_i)_{i=1}^n$.
The reverse inclusion follows from the original identity
$\beta\alpha=0$, so the localized sequence is exact.
The local exactness criterion proves the original sequence exact.
If the covering list is empty, the unit-ideal hypothesis gives $R=0$,
and every unital $R$-module is zero, so the same assertion holds.
\end{proof}